\chapter[(Аштавакра)]{Настанови для пізнання себе (Аштавакра)}

[1.2] Щоб стати цілковито вільним, уникай полону чуттєвих насолод, наче отрути. Плекай терпимість, щирість, співчуття, простоту й чесність.

[1.3] Ти --- не земля, не вода, не вогонь і не повітря. І навіть не порожній простір. Звільнення --- це пізнання себе як чистої Свідомості, що є свідком усього цього.

[1.4] Якщо ти зумієш усвідомити себе чистою Свідомістю, переставши ототожнювати себе з тілом, то вмить досягнеш щастя, спокою та звільнення.

[1.5] Ти не належиш до жодної касти й не обмежений жодними обов'язками. Твоя природа невидима для очей, неприв'язана та безформна. Вона --- свідок усього. Будь щасливий.

[1.6] Правильне й неправильне, насолода та біль існують лише в умі й не стосуються тебе. Не ти дієш і не ти зазнаєш наслідків дій, тож ти завжди вільний.

[1.7] Ти --- Єдиний Свідок і завжди вільний. Причина твоєї неволі --- в тому, що ти вважаєш себе чимось іншим.

[1.8] Думка «я дію» --- мов укус отруйної змії. Усвідомлення «я нічого не роблю» --- мудрість-протиотрута. Будь щасливий.

[1.9] Самого розуміння «я --- Єдина Свідомість» досить, щоб у вогні однієї миті згоріли невігластво й страждання. Будь щасливий.

[1.10] Ти --- Безмежна Свідомість, блаженство, найвище захоплення. У тобі Всесвіт набуває форми, немов змія, яку помилково вбачають у мотузці. Будь щасливий.

[1.11] Недарма кажуть: «Думка формує твою дійсність». Якщо вважаєш себе обмеженим, то обмежений. Якщо вважаєш себе безмежним, то безмежний.

[1.12] Твоя справжня природа --- досконала, вільна, всюдисуща й єдина, ні з чим не ототожнена, без бажань і сповнена спокою. Самсара --- ілюзія, породжена невіглаством. Це існування без пізнання себе.

[1.13] Споглядай себе як саму лише Свідомість --- єдність як таку. Відмовся від уявлення, ніби ти відокремлений, ніби ти --- свідомість, що блукає. Відмовся від думки, що існує «всередині» та «зовні».

[1.14] Ти давно перебуваєш у пастці ототожнення з тілом, думаючи: «Я --- людина». Нехай знання «я --- лише Свідомість» стане мечем, який тебе звільнить. Будь щасливий.

[1.15] Ти завжди був вільним, світлосяйним і незаплямованим. Намагання силоміць заспокоїти розум --- саме те, що тримає тебе в неволі.

[1.16] Ти --- Свідомість, сутність Всесвіту. Дійсність існує завдяки тобі. Не обмежуй свого бачення.

[1.17] Ти нічим не зумовлений, незмінний, безформний, нерухомий, незбагненний і спокійний. Нічого не бажай. Просто будь чистою Свідомістю.

[1.18] Форма видається справжньою, тоді як безформне є вічним. Усвідомивши це, ти більше не потрапиш в оману.

[1.19] Як у дзеркалі є образ, а перед ним --- те, що віддзеркалюється, так Найвище Я існує всередині й поза межами тіла.

[1.20] Як усюдисущий простір існує і всередині глечика, і поза ним, так позачасова Єдність охоплює все.
